% ECONOMICS_PROBLEM: {"id": "J65-312", "title": "Optimal unemployment and social insurance under heterogeneity", "jel_code": "J65", "source_id": "OP-0045"}
% !TeX program = xelatex
\documentclass[11pt,a4paper]{article}
\usepackage[margin=23mm]{geometry}
\usepackage{fontspec}
\setmainfont{Latin Modern Roman}
\usepackage{amsmath,amssymb,mathtools}
\usepackage{xcolor,enumitem,booktabs,longtable,array,xurl,hyperref}
\definecolor{muted}{HTML}{53636C}
\hypersetup{colorlinks=true,urlcolor=blue,linkcolor=blue}
\setlength{\parindent}{0pt}
\setlength{\parskip}{5pt}
\newcommand{\R}{\mathbb R}
\newcommand{\E}{\mathbb E}
\newcommand{\N}{\mathbb N}
\newcommand{\ind}{\mathbf 1}
\newcommand{\Var}{\operatorname{Var}}
\newcommand{\Cov}{\operatorname{Cov}}
\newcommand{\supp}{\operatorname{supp}}
\DeclareMathOperator*{\argmin}{arg\,min}
\DeclareMathOperator*{\argmax}{arg\,max}
\newcommand{\entrymeta}[3]{{\sffamily\small\color{muted}\textbf{\texttt{#1}}\quad|\quad #2\par #3\par}}
\newcommand{\Problemref}[1]{\texttt{#1}}
\newcommand{\JELref}[2]{\texttt{#2} (JEL \texttt{#1})}
\begin{document}
\section*{Optimal unemployment and social insurance under heterogeneity}
\textbf{Problem ID:} \texttt{J65-312}\quad \textbf{JEL:} \texttt{J65}\par
% BEGIN ECONOMICS STATEMENT
\label{problem:OP-0045}
\label{atlas:prob:P5}
\noindent\textbf{Original atlas identifier:} P5; entry 45. \textbf{Field:} Public finance and taxation. \textbf{Original tractability label:} Medium.

\noindent\textbf{Classification:} Parameterized theoretical/empirical research agenda; not a certified unsolved theorem.

\subsection*{Definitions and assumptions}
Let an unemployed worker have observable history $x_t$, unobserved type $\theta$, assets $a_t$, unemployment duration $d_t$, and aggregate state $z_t$ at date $t$. A benefit policy $b(x_t,d_t,z_t)$ is measurable only with respect to information legally available to the insurer. A model $m$ specifies liquidity constraints, search effort, acceptance decisions, employment hazards, matching technology, wage setting, and equilibrium selection. Fix discount factor $\beta\in(0,1)$, nonnegative type weights $w(\theta)$, cardinal utility $u(c,\ell;\theta)$ over consumption $c$ and leisure $\ell$, and public-fund value $\lambda$ in utility units per currency unit. Let $G_t(b;m)$ denote net public outlays, including benefit payments, administration, and induced tax-receipt changes. Assume expected discounted utilities and outlays are finite.

\subsection*{Formal research question}
\emph{Editorial formulation: the mathematical target below is not a theorem quoted from the references.}

Among legally feasible benefit schedules financed by an explicitly specified payroll tax or public budget, characterize
\[
 \sup_{b\in\mathcal B}\mathbb E_m\sum_{t=0}^{\infty}\beta^t
 \left[w(\theta)u(c_t,\ell_t;\theta)-\lambda G_t(b;m)\right].
\]
Identify the welfare gain from allowing dependence on $x_t$, $d_t$, and $z_t$ separately. Determine what variation separates liquidity effects from substitution in search and from changes in market tightness. Specify an identified model set $\mathcal M$ when benefit reforms recover only local average responses; compare benefit schedules over that set and report which policy rankings are robust. Do not infer an optimal individualized rule from a duration elasticity without the consumption, welfare, and equilibrium primitives that make it sufficient.

\subsection*{Interpretation and scope}
Personalized insurance can be implemented only using admissible information; an optimum conditioning on an unobserved type is an unattainable benchmark unless incentive-compatible disclosure is included. Changes in job finding can reflect search effort, job availability, wage responses, or acceptance behavior, with different welfare implications. Business-cycle policy therefore requires a financing and wage-setting closure, not merely multiplication of a partial-equilibrium formula by unemployment. Administrative heterogeneity must not be interpreted as preference heterogeneity without measurement and selection assumptions. Duration reforms can have anticipation and exhaustion effects, which need a stated dynamic design. The intended extension is feasible heterogeneity combined with aggregate labor-demand conditions, while existing insurance/search and matching-model benchmarks remain recognized solved cases.

\subsection*{Source audit and status}
All original sources are real. [S2] has a publisher erratum and should be consulted with its correction. The 2018 author-year in [S3] is disambiguated to the Theory paper, distinct from its Applications companion. These antecedents do not establish an unresolved general theorem in 2026.

\subsection*{References}
\begin{enumerate}[label={[S\arabic*]},leftmargin=*,itemsep=0.6em]
\item Martin Neil Baily (1978). Some Aspects of Optimal Unemployment Insurance. Journal of Public Economics 10(3), 379--402. DOI: 10.1016/0047-2727(78)90053-1.
\url{https://www.sciencedirect.com/science/article/pii/0047272778900531}
\newline\emph{Locator and support limits:} Article: insurance/search benchmark; worker heterogeneity and macroeconomic equilibrium require extensions.
\item Raj Chetty (2008). Moral Hazard versus Liquidity and Optimal Unemployment Insurance. Journal of Political Economy 116(2), 173--234. DOI: 10.1086/588585. Erratum: 116(6), DOI 10.1086/597025.
\url{https://www.journals.uchicago.edu/doi/10.1086/588585}
\newline\emph{Locator and support limits:} Abstract and publisher correction notice: liquidity versus search responses; consult the erratum before reproducing numerical formulas.
\item Camille Landais, Pascal Michaillat, and Emmanuel Saez (2018). A Macroeconomic Approach to Optimal Unemployment Insurance: Theory. American Economic Journal: Economic Policy 10(2), 152--181. DOI: 10.1257/pol.20150088.
\url{https://www.aeaweb.org/articles?id=10.1257/pol.20150088}
\newline\emph{Locator and support limits:} Abstract: matching-model tightness correction. The original author-year also fits the companion Applications article; Theory is the explicit selection.
\end{enumerate}

\subsection*{Common conventions: Source mathematical conventions}

Unless an entry states otherwise, agent and firm sets are finite. State and action spaces are measurable, strategies and outcome maps are measurable, probability laws are specified, and expectations exist for the targets under discussion. Infinite-horizon discount factors lie in $(0,1)$ and discounted objectives are finite. Norms, tolerances and complexity input sizes must be specified for computational comparisons. Constants or primitives that appear locally are inputs, not empirical facts asserted by this catalogue.

\textbf{Identification.} A statistical model $m\in\mathcal M$ implies an observable law $P_m$ and target $\theta(m)$. At law $P$, its identified set is
\[
\Theta(P;\mathcal M)=\{\theta(m):m\in\mathcal M,\ P_m=P\}.
\]
Point identification holds when this set is a singleton, or equivalently when $P_{m_1}=P_{m_2}$ implies $\theta(m_1)=\theta(m_2)$ for all compatible models. Sharp bounds mean bounds attained, or approached when the boundary is not attained, by compatible models. Writing this set is a definition, not a solution: the substantive task is to establish informative restrictions and derive or estimate the set. An unrestricted-confounding model may give only trivial bounds.

\textbf{Inference.} Identification, estimability and valid uncertainty quantification are separate questions. Fix $\alpha\in(0,1)$. A confidence procedure $C_n$ over a declared law class $\mathcal P$ has asymptotic uniform coverage at level $1-\alpha$ when
\[
\liminf_{n\to\infty}\inf_{P\in\mathcal P}\Pr_P\{\theta(P)\in C_n\}\geq1-\alpha.
\]
For set-identified targets the coverage event must be defined separately, for example coverage of the full identified set. If an entry uses identification-indexed models instead of uniquely indexed laws, the same coverage convention is applied over those models. Pointwise limits do not establish uniform validity, and asymptotic validity does not guarantee finite-sample precision. Sample sizes, dependence structures and weak-identification sequences are part of each application.

\textbf{Counterfactuals.} $Y_i(a)$ is a potential outcome under a specified intervention, including its timing, implementation and versions. It is not $Y_i$ conditional on voluntarily choosing $a$. A full intervention vector permits interference; shorthand scalar treatment notation does not silently impose no interference. Assignment, selection, attrition, anticipation and measurement must be specified. A claim about an instrument's local effect is not automatically a population policy effect.

\textbf{Economic equilibrium.} A counterfactual model includes best responses, constraints, feasibility, market clearing, state transitions and expectations consistent with the stated information. When several equilibria exist, either a declared selector is an input or the target is set-valued. Comparisons across policy regimes must use the stated selection convention. Reduced-form relationships are not presumed invariant after an equilibrium-changing intervention.

\textbf{Optimization and welfare.} Optimization is over the stated feasible set. If attainment is not established, $\sup$ or $\inf$ and $\varepsilon$-optimal choices replace an asserted maximizer or minimizer. For example, with value $V=\sup_{a\in\mathcal A}W(a)<\infty$, an $\varepsilon$-optimal set is $\{a\in\mathcal A:W(a)\geq V-\varepsilon\}$ for $\varepsilon>0$. Benefits, costs, risk and health or environmental outcomes can be aggregated only using declared units and conversion weights. Interpersonal utility weights, the treatment of fiscal transfers and foreign profits, discounting, and the behavioral welfare criterion are explicit inputs. A comparative-static derivative additionally requires existence and appropriate differentiability of the selected counterfactual.

\textbf{Sources.} Local labels [S1], [S2], and so forth restart in every question. Locators identify the relevant abstract, model, section or result at the level actually checked; a bibliographic or abstract-level verification is not represented as inspection of an entire proof. Working papers are distinguished from later publications and changing versions. Source summaries are paraphrased. All URL checks and editorial formulations refer to the audit date stated above.
% END ECONOMICS STATEMENT
\end{document}
